\begin{tikzpicture}[
  font=\sffamily,
  cell/.style={
    rectangle, rounded corners=3pt, line width=0.9pt,
    text width=6.15cm, align=left, inner sep=7pt, minimum height=2.25cm, font=\sffamily\scriptsize
  }
]
  % Etiquetas de ejes
  \node[font=\sffamily\bfseries\footnotesize, ceindigo] at (3.4, 3.65) {FACTORES NEGATIVOS (Riesgos / Frenos)};
  \node[font=\sffamily\bfseries\footnotesize, ceteal!90!black] at (10.1, 3.65) {FACTORES POSITIVOS (Ventajas / Palancas)};
  \node[font=\sffamily\bfseries\footnotesize, ceindigo, rotate=90] at (-0.65, 2.0) {ORIGEN INTERNO};
  \node[font=\sffamily\bfseries\footnotesize, cecrimson, rotate=90] at (-0.65, -0.75) {ORIGEN EXTERNO};

  % Fila Interna: Debilidades -> Corregir | Fortalezas -> Mantener
  \node[cell, draw=cecrimson, fill=cepanelcrimson] (d) at (3.4, 2.0) {
    \textcolor{cecrimson}{\bfseries DEBILIDADES (Internas)} $\longrightarrow$ \textbf{CORREGIR}\\[3pt]
    \color{ceslate} Carencias del equipo, deuda técnica, falta de marca o liquidez limitada.\\[2pt]
    \textit{Estrategia CAME de Supervivencia / Reorientación.}
  };
  \node[cell, draw=ceteal, fill=cepanelteal] (f) at (10.1, 2.0) {
    \textcolor{ceteal!90!black}{\bfseries FORTALEZAS (Internas)} $\longrightarrow$ \textbf{MANTENER}\\[3pt]
    \color{ceslate} Patentes, equipo técnico mixto, costes bajos, comunidad fiel o red propia.\\[2pt]
    \textit{Estrategia CAME Defensiva u Ofensiva.}
  };

  % Fila Externa: Amenazas -> Afrontar | Oportunidades -> Explotar
  \node[cell, draw=ceamber, fill=cepanelamber] (a) at (3.4, -0.75) {
    \textcolor{ceamber!90!black}{\bfseries AMENAZAS (Externas)} $\longrightarrow$ \textbf{AFRONTAR}\\[3pt]
    \color{ceslate} Entrada de gigantes, guerras de precios, cambios regulatorios o recesión.\\[2pt]
    \textit{Transformar riesgos mediante pivote o nicho.}
  };
  \node[cell, draw=ceindigo, fill=cepanel] (o) at (10.1, -0.75) {
    \textcolor{ceindigo}{\bfseries OPORTUNIDADES (Externas)} $\longrightarrow$ \textbf{EXPLOTAR}\\[3pt]
    \color{ceslate} Nuevas leyes, crecimiento demográfico, brechas tecnológicas o caída de un rival.\\[2pt]
    \textit{Ataque directo apoyado en Fortalezas (Estrategia Ofensiva).}
  };
\end{tikzpicture}